% ---------------------------------------------------------------------
% HAND-TRANSCRIBED SOURCE PAGE — scan idx 71 (printed p. 267)
%
% Hand-transcribed from the leaf at 500 dpi.
% Running head "FAISCEAUX ALGEBRIQUES COHERENTS   267" is page furniture, excluded.
%
% NOTATION: two decorative faces occur and are kept apart.
%   - the sheaf-from-module functor prints as the same blackletter cap used for A/B/C on the
%     earlier pages of this workspace -> \mathfrak{A}, matching \mathfrak{C} in "M^* \in C"
%   - the structure sheaf prints as a ROUND script O, not the angular blackletter O -> \mathcal{O}
% The distinction is deliberate; do not unify them.
%
% Inequalities are the double-bar \geqq of this journal.
%
% PRINTING ARTIFACT, not transcribed: in the first display there is a small speck of ink above
% the N of "\mathfrak{A}(N)". It is not a diacritic — nothing in the surrounding text carries one,
% and the same expression is set clean two displays later.
% ---------------------------------------------------------------------
Comme les propriétés (i) à (iv) caractérisent les foncteurs dérivés du foncteur $\mathrm{Hom}_S(M, \Omega)$, on a $E^q(M) \approx \mathrm{Ext}_S^q(M, \Omega)$, ce qui démontre le Théorème.

\textsc{Corollaire} 1. \textit{Si $M$ vérifie} (TF), $H^q(M)$ \textit{est isomorphe à l'espace vectoriel dual de} $\mathrm{Ext}_S^{r-q}(M, \Omega)_0$ \textit{pour tout} $q \geqq 1$.

En effet, nous savons que $H^q(M)$ est un espace vectoriel de dimension finie dont le dual est isomorphe à $\mathrm{Ext}_S^{r-q}(M, \Omega)_0$ .

\textsc{Corollaire} 2. \textit{Si $M$ vérifie} (TF), \textit{les} $T^q(M)$ \textit{sont des $S$-modules gradués de type fini pour} $q \geqq 1$, \textit{et} $T^0(M)$ \textit{vérifie} (TF).

On peut remplacer $M$ par un module de type fini sans changer les $B^q(M)$, donc les $T^q(M)$. Les $\mathrm{Ext}_S^{r-q}(M, \Omega)$ sont alors des $S$-modules de type fini, et l'on a $M^* \in \mathfrak{C}$, d'où le Corollaire.

\begin{center}
\textbf{\S5. Applications aux faisceaux algébriques cohérents}
\end{center}

\textbf{73. Relations entre les foncteurs $\mathrm{Ext}_S^q$ et $\mathrm{Ext}_{\mathcal{O}_x}^q$ .} Soient $M$ et $N$ deux $S$-modules gradués. Si $x$ est un point de $X = \mathbf{P}_r(K)$, nous avons défini au n\textsuperscript{o} 57 les $\mathcal{O}_x$-modules $M_x$ et $N_x$ ; nous allons mettre en relation les $\mathrm{Ext}_{\mathcal{O}_x}^q(M_x , N_x)$ avec le $S$-module gradué $\mathrm{Ext}_S^q(M, N)$:

\textsc{Proposition} 1. \textit{Supposons que $M$ soit de type fini. Alors:}

(a) \textit{Le faisceau} $\mathfrak{A}(\mathrm{Hom}_S(M, N))$ \textit{est isomorphe au faisceau} $\mathrm{Hom}_{\mathcal{O}}(\mathfrak{A}(M), \mathfrak{A}(N))$.

(b) \textit{Pour tout} $x \in X$, \textit{le $\mathcal{O}_x$-module} $\mathrm{Ext}_S^q(M, N)_x$ \textit{est isomorphe au $\mathcal{O}_x$-module} $\mathrm{Ext}_{\mathcal{O}_x}^q(M_x , N_x)$.

Définissons d'abord un homomorphisme $\iota_x: \mathrm{Hom}_S(M, N)_x \rightarrow \mathrm{Hom}_{\mathcal{O}_x}(M_x , N_x)$. Un élément du premier module est une fraction $\varphi/P$, avec $\varphi \in \mathrm{Hom}_S(M, N)_n$ , $P \in S(x)$, $P$ homogène de degré $n$; si $m/P'$ est un élément de $M_x$ , $\varphi(m)/PP'$ est un élément de $N_x$ qui ne dépend que de $\varphi/P$ et de $m/P'$, et l'application $m/P' \rightarrow \varphi(m)/PP'$ est un homomorphisme $\iota_x(\varphi/P): M_x \rightarrow N_x$ ; ceci définit $\iota_x$ . D'après la Proposition 5 du n\textsuperscript{o} 14, $\mathrm{Hom}_{\mathcal{O}_x}(M_x , N_x)$ peut être identifié à

\begin{equation*}
\mathrm{Hom}_{\mathcal{O}}(\mathfrak{A}(M), \mathfrak{A}(N))_x \ ;
\end{equation*}

cette identification transforme $\iota_x$ en

\begin{equation*}
\iota_x : \mathfrak{A}(\mathrm{Hom}_S(M, N))_x \longrightarrow \mathrm{Hom}_{\mathcal{O}}(\mathfrak{A}(M), \mathfrak{A}(N))_x \ ,
\end{equation*}

et l'on vérifie facilement que la collection des $\iota_x$ est un homomorphisme

\begin{equation*}
\iota : \mathfrak{A}(\mathrm{Hom}_S(M, N)) \longrightarrow \mathrm{Hom}_{\mathcal{O}}(\mathfrak{A}(M), \mathfrak{A}(N)).
\end{equation*}

Lorsque $M$ est un module libre de type fini, $\iota_x$ est bijectif: en effet, il suffit de le voir lorsque $M = S(n)$, auquel cas c'est immédiat.

Si maintenant $M$ est un $S$-module gradué de type fini quelconque, choisissons une résolution de $M$:

\begin{equation*}
\cdots \rightarrow L^{q+1} \rightarrow L^q \rightarrow \cdots \rightarrow L^0 \rightarrow M \rightarrow 0,
\end{equation*}

où les $L^q$ soient libres de type fini, et considérons le complexe $C$ formé par les $\mathrm{Hom}_S(L^q, N)$. Les groupes de cohomologie de $C$ sont les $\mathrm{Ext}_S^q(M, N)$; autrement
